\section{Introduction}
\label{sec:introduction}

Mass spectrometry imaging (MSI) records a mass spectrum at each sampled tissue
coordinate, allowing molecular intensity patterns to be examined without
preselecting a single target molecule~\citep{caprioli1997maldi,schwamborn2010molecular}.
The resulting data are both spatial and high-dimensional: each tissue section
contains many measured locations, and each location contains intensities for
hundreds to tens of thousands of mass-to-charge ($m/z$) bins. A downstream
analysis therefore requires a defensible way to reduce the spectral dimension
without discarding weak but informative signals.

Classical peak-picking methods identify local spectral maxima using processing
choices such as smoothing, signal-to-noise rules and thresholds. Learning-based
alternatives instead optimise a representation of the spectra and use the
trained model to rank input bins. The msiPL method applies a variational
autoencoder (VAE) to individual spectra and derives a peak ranking from network
weights~\citep{abdelmoula2021peak}. S3PL supplies a spatial-spectral patch to a
three-dimensional convolutional model~\citep{weigand2026spatial}. These methods
motivate two related but distinct questions: whether local neighbourhood
context improves unsupervised peak selection, and whether a nonlinear
explanation of a learned spectral representation yields a more useful ranking
than a first-layer weight shortcut.

This study investigates those questions with centre-only and contextual VAE
models on glioblastoma (GBM) and colorectal adenocarcinoma (CAC) MSI sections.
The contextual model reconstructs each central spectrum after receiving both
the centre and a summary of neighbouring spectra. Fixed means, larger windows,
zero and shuffled controls, and learned attention are used to test what the
neighbourhood contributes. Separately, a two-component Gaussian mixture model
(GMM) is fitted to frozen latent means without expert labels. Integrated
Gradients (IG) then attributes the GMM posterior through the complete encoder
to the original spectral bins~\citep{sundararajan2017axiomatic}. Expert masks
are introduced only during final evaluation.

The central research question is whether neighbourhood context provides a
repeatable benefit for unsupervised MSI peak selection. Two supporting
questions ask whether nonlinear GMM-targeted attribution improves the evaluated
pipeline over reproduced weight-based peak learning, and how spatial controls,
training instability and measured workflow cost qualify the comparison.

The report makes four bounded contributions. First, it provides matched
centre-only and contextual VAE comparisons across eight CAC and eight GBM
sections. Second, it evaluates a label-free GMM/IG ranking pipeline against
legacy msiPL and a first-layer L2 comparator at matched peak budgets. Third, it
uses window, zero-context, shuffled-context, learned-attention and S3PL controls
to separate model sensitivity from evidence of a beneficial local topology.
Fourth, it records baseline stability and native workflow resource measurements
so that peak-quality claims are not separated from their computational and
reproducibility limits.

The evaluated VAE/GMM/IG package improves matched-budget spatial peak scores
over reproduced legacy msiPL on both collections. However, the benefit that can
be attributed specifically to local neighbourhood structure is not consistent
across CAC and GBM. This mixed result is treated as the main scientific finding,
not as a failed attempt to obtain a positive result. The remainder of the report
reviews the relevant methods, defines the controlled experiments and proposal
changes, presents section- and patient-aware evidence, and states the limits of
the biological and computational conclusions.

